\section{Packet fixed-point series and their dyadic poles}
\label{sec:packet-fixed-point-series}

The affine coordinates retain useful analytic information when summed over
all words of a prescribed length and exponent sum.  The total affine offset
has a rational generating function, whereas the total fixed point has
infinitely many poles.  We construct both functions in the original packet
coordinates and determine their exact domains of power-series convergence.
The rational recursion and formal dilation equation occur in the local
\emph{Affine Packets} note, Sections~15--17, and in raw Chatnotes
\emph{Implications of Affine Formula}, lines~432--475.  The continuation
below extends the note's one-letter pole calculation to every packet length.
These statements concern sums over rational fixed points, with repeated
words retained; they are not sums over positive integral cycles alone.

\subsection{The rational offset family}

For $m\ge1$ and $A\ge m$, retain $C(p)$, $D=2^A-3^m$, and
$x_p=C(p)/D$ from \eqref{eq:packet-C-D-x}.  Put
\[
 c(p)=2^{-A}C(p),\qquad
 G_m(A)=|\Pi_{m,A}|,\qquad
 S_m(A)=\sum_{p\in\Pi_{m,A}}c(p),\qquad
 X_m(A)=\sum_{p\in\Pi_{m,A}}x_p.
\]
Set these three coefficients to zero for $A<m$.  Their generating
functions are $G_m(q)=\sum_A G_m(A)q^A$, $F_m(q)=\sum_A S_m(A)q^A$,
and $X_m(q)=\sum_A X_m(A)q^A$.  The same letter $X_m$ with an integer
argument denotes a coefficient; with $q$ it denotes the series.  Define
\[
 u(q)=\frac q{1-q},\qquad v(q)=\frac q{2-q}.
\]

\begin{proposition}[Packet sums and their two-state representation]
\label{prop:packet-total-offset-series}
As formal series, and hence as rational functions where defined,
\begin{align}
 G_m(q)&=u(q)^m,\label{eq:analytic-packet-cardinality}\\
 F_{m+1}(q)&=v(q)(3F_m(q)+G_m(q)),\qquad F_1(q)=v(q),
 \label{eq:analytic-packet-recursion}\\
 F_m(q)&=\sum_{k=0}^{m-1}3^{m-1-k}u(q)^k v(q)^{m-k}.
 \label{eq:analytic-packet-closed}
\end{align}
In particular, with column vectors,
\[
 \binom{F_{m+1}}{G_{m+1}}
 =\begin{pmatrix}3v&v\\0&u\end{pmatrix}\binom{F_m}{G_m},
 \qquad \binom{F_1}{G_1}=\binom vu.
\]
The bivariate formal series is
\begin{equation}\label{eq:analytic-bivariate-offset}
 \sum_{m\ge1}F_m(q)z^{m-1}
 =\frac{v(q)}{(1-3v(q)z)(1-u(q)z)}.
\end{equation}
For $m\ge2$, $F_m$ has exactly two poles: order $m-1$ at $1$ and
order $m$ at $2$.  For $m=1$ its only pole is the simple pole at $2$.
At the regular point $q=1/2$, $F_m(1/2)=m/3$.
\end{proposition}
\begin{proof}
A word is a composition of $A$ into $m$ positive parts, so
$G_m(A)=\binom{A-1}{m-1}$ and its series is $u^m$.  Appending $a\ge1$
to a word $p$ gives $c(p,a)=2^{-a}(3c(p)+1)$.  Sum first over
$p\in\Pi_{m,A-a}$ and then over $1\le a\le A-m$.  Each coefficient
is a finite sum, and multiplication by $\sum_{a\ge1}2^{-a}q^a=v(q)$
gives \eqref{eq:analytic-packet-recursion}.  Induction proves
\eqref{eq:analytic-packet-closed}.  Summing that finite-sum formula over
$m$ gives the product of the two geometric series in
\eqref{eq:analytic-bivariate-offset}; both denominators have constant
term one in $z$.

At $q=1$, only the $k=m-1$ term of \eqref{eq:analytic-packet-closed}
has pole order $m-1$, and its coefficient relative to
$(1-q)^{-(m-1)}$ is one.  At $q=2$, only the $k=0$ term has order
$m$, and its coefficient relative to $(2-q)^{-m}$ is
$3^{m-1}2^m$.  Both coefficients are nonzero.  The finite-sum formula
has no other denominator zeros.
Finally $u(1/2)=1$ and $v(1/2)=1/3$, so every one of the $m$ terms
in the finite sum equals $1/3$.
\end{proof}

\subsection{Formal inversion and convergence domains}

\begin{proposition}[The exact dilation map]\label{prop:packet-dilation-inverse}
For each $m\ge1$, the map
\[
 \mathcal D_m:q^m\mathbb C[[q]]\longrightarrow q^m\mathbb C[[q]],
 \qquad f(q)\longmapsto f(2q)-3^m f(q)
\]
is a continuous complex-linear bijection for the $q$-adic topology.
Its inverse is
\[
 \mathcal D_m^{-1}\!\left(\sum_{A\ge m}h_Aq^A\right)
 =\sum_{A\ge m}\frac{h_A}{2^A-3^m}q^A.
\]
Both maps preserve the order of every nonzero formal series.  Their
kernels are zero and every fibre is a singleton.  In particular,
\begin{equation}\label{eq:packet-dilation-equation}
 X_m(2q)-3^mX_m(q)=F_m(2q).
\end{equation}
The radius of convergence of $X_1$ is $2$.  For every $m\ge2$ the
radius of $X_m$ is exactly $1$.  Consequently
\eqref{eq:packet-dilation-equation}, interpreted as an equality of the
defining convergent power series, holds on $|q|<1$ for $m=1$, and on
$|q|<1/2$ for $m\ge2$.
\end{proposition}
\begin{proof}
The coefficient multiplier $2^A-3^m$ is nonzero: its first term is even
and its second odd.  The two displayed maps are inverse coefficient by
coefficient and preserve the least index of a nonzero coefficient.
This proves all the formal topological and linear assertions.  It does
not assert preservation of multiplication.  The identity
\[
 X_m(A)=\frac{S_m(A)}{1-3^m2^{-A}}
\]
then gives \eqref{eq:packet-dilation-equation}.

For $m=1$, $X_1(A)=(2^A-3)^{-1}$ and
$|X_1(A)|^{1/A}\to1/2$.  For fixed $m$, every word satisfies
\[
 0<c(p)\le \sum_{j=0}^{m-1}3^{m-1-j}2^{-(m-j)}
 =(3/2)^m-1.
\]
Here $A-S_j\ge m-j$ because the remaining letters are positive.
For all sufficiently large $A$, $1-3^m2^{-A}\ge1/2$.  Therefore
\[
 0<X_m(A)\le 2\bigl((3/2)^m-1\bigr)\binom{A-1}{m-1}.
\]
The finitely many earlier coefficients do not affect the radius.  This
polynomial upper bound gives convergence for $|q|<1$.  For $m\ge2$
and each $A\ge m$, there is a word whose last letter is $1$, for example
$(A-m+1,1,\ldots,1)$.  Its final numerator summand gives $c(p)\ge1/2$.
Once $2^A>3^m$, all packet terms are positive, so $X_m(A)\ge1/2$.
The terms therefore cannot tend to zero when $|q|>1$.  The radius is
exactly one, and the claimed domains follow after replacing $q$ by $2q$.
The same lower bound shows failure of the term test at every point of
the boundary circle for $m\ge2$.  For $m=1$, the terms on $|q|=2$
have magnitudes $2^A/|2^A-3|\to1$, so they also fail that test.
\end{proof}

For example,
\[
 X_2(A)=\frac{2^A+3A-5}{2^A-9}\quad(A\ge2),
\]
because $C(a,A-a)=3+2^a$.  Thus $X_2(A)\to1$.  At $q=3/4$, the
terms defining $X_2(2q)$ do not tend to zero.  This disproves the
whole-unit-disc convergence assertion appended to the formal identity
in the local note's Theorem~17.1.  The formal identity itself is exact.

\subsection{Meromorphic continuation with exact pole orders}

\begin{theorem}[The dyadic pole family]\label{thm:packet-dyadic-poles}
For every $m\ge1$, $X_m$ has a meromorphic continuation to the complex
plane.  When $m=1$ its poles are precisely $2^j$, $j\ge1$, all simple.
When $m\ge2$, it has a pole of order $m-1$ at $1$ and a pole of order
$m$ at each $2^j$, $j\ge1$, with no other poles.  More precisely,
\begin{align}
 \lim_{q\to1}(1-q)^{m-1}X_m(q)&=1 &&(m\ge2),
 \label{eq:packet-pole-one}\\
 \lim_{q\to2^j}(2^j-q)^mX_m(q)&=3^{mj-1}2^{jm}
 &&(m\ge1,\ j\ge1).
 \label{eq:packet-dyadic-leading}
\end{align}
In particular, none of the functions $X_m$ is rational.
\end{theorem}
\begin{proof}
Choose an integer $L\ge m$ with $2^L>3^m$; for example $L=2m$.
Set
\[
 P_m(q)=\sum_{A=m}^{L-1}X_m(A)q^A,\qquad
 H_m(q)=F_m(q)-\sum_{A=m}^{L-1}S_m(A)q^A.
\]
The empty sums are zero.  The rational function $H_m$ vanishes to
order at least $L$ at zero and has exactly the same poles and principal
parts as $F_m$.  The following explicit expression continues $X_m$:
\begin{equation}\label{eq:packet-dyadic-continuation}
 X_m(q)=P_m(q)+\sum_{r\ge0}3^{mr}H_m(q/2^r).
\end{equation}
Indeed, for $A\ge L$,
\[
 \frac1{1-3^m2^{-A}}=\sum_{r\ge0}3^{mr}2^{-rA}.
\]
For $|q|<1$, the sum of absolute values of this double series is at most
\[
 \frac{1}{1-3^m2^{-L}}\sum_{A\ge L}S_m(A)|q|^A<\infty.
\]
The previous proposition supplies the polynomial bound needed here.
Interchanging the sums proves \eqref{eq:packet-dyadic-continuation}
near zero.

For an explicit normal-convergence estimate put
\[
 B_{m,L}=2^L H_m(1/2)>0,\qquad \eta=3^m2^{-L}<1.
\]
All coefficients of $H_m$ from degree $L$ onward are positive.  Hence
$|H_m(z)|\le B_{m,L}|z|^L$ for $|z|\le1/2$.
For $|q|\le R$ and any integer $N\ge0$ with $R/2^N\le1/2$,
\begin{equation}\label{eq:packet-continuation-error}
 \sup_{|q|\le R}\left|\sum_{r=N}^{\infty}3^{mr}H_m(q/2^r)\right|
 \le \frac{B_{m,L}R^L\eta^N}{1-\eta}.
\end{equation}
The constant is explicitly rational:
$B_{m,L}=2^L(m/3-\sum_{A=m}^{L-1}S_m(A)2^{-A})$.
This is a summable geometric bound uniform on every compact set.
Only finitely many earlier summands remain.  Each is rational, with
poles confined to the positive dyadic points.  Away from those points,
the sum is locally uniformly convergent and holomorphic.  Near each
such point, all but finitely many summands are holomorphic and their
tail converges uniformly, so the sum is meromorphic there.

At $q=1$ only $r=0$ can be singular, giving
\eqref{eq:packet-pole-one}.  At $q=2^j$, $j\ge1$, the term $r=j-1$
inherits the pole of $F_m$ at $2$, of order $m$.  Its leading
coefficient in powers of $(2^j-q)^{-1}$ is
\[
 3^{m(j-1)}3^{m-1}(2^j)^m=3^{mj-1}2^{jm}.
\]
For $m\ge2$, the only other singular term is $r=j$, whose order is
$m-1$; it cannot cancel the leading pole.  For $m=1$ that term is
regular.  No other terms are singular.  This proves the exact pole
list and both limits.  A rational function has only finitely many
poles, whereas this list is infinite.
\end{proof}

The same expression also proves that the dilation equation holds globally
as an identity of meromorphic functions, by analytic continuation from
its disc of convergent series.  This is a different assertion from
convergence of the defining series beyond its radius.  For $m=1$, taking
$L=2$ gives $P_1(q)=-q$ and recovers the source's formula
\[
 X_1(q)=-q+\sum_{r\ge0}3^r
 \frac{(q/2^{r+1})^2}{1-q/2^{r+1}}.
\]
Thus the one-letter pole calculation and the general packet continuation
are related by an explicit specialization, without changing either series.
